% Generated draft: requires mathematical and source review.
\documentclass[11pt]{article}
\usepackage[margin=25mm]{geometry}
\usepackage{fontspec}
\setmainfont{Latin Modern Roman}
\usepackage{amsmath,amssymb,mathtools,unicode-math}
\setmathfont{Latin Modern Math}
\usepackage{enumitem,needspace,xurl,hyperref}
\setlength{\parindent}{0pt}
\setlength{\parskip}{5pt}
\setlength{\emergencystretch}{3em}
\newcommand{\N}{\mathbb{N}}
\newcommand{\Z}{\mathbb{Z}}
\newcommand{\Q}{\mathbb{Q}}
\newcommand{\R}{\mathbb{R}}
\newcommand{\C}{\mathbb{C}}
\newcommand{\F}{\mathbb{F}}
\newcommand{\A}{\mathbb{A}}
\newcommand{\PP}{\mathbb{P}}
\newcommand{\E}{\mathbb{E}}
\newcommand{\eps}{\varepsilon}
\newcommand{\dd}{\,\mathrm{d}}
\begin{document}
\section*{UnsolvedMath ID: 30003068}
\textbf{Source problem code:} catalogue entry
\subsection{Definitions and mathematical statement}
The supplied record asks the following two nonemptiness questions in the setting of PBW structures. Let \(\mathfrak g\) be a simple finite-dimensional Lie algebra in the intended source setting, and let \(S_{\mathrm{gm}}\) denote its source-defined set of global monomial degrees.
\begin{enumerate}
\item For every such \(\mathfrak g\), is
\[
S_{\mathrm{gm}}\ne\varnothing?
\]
\item For every simply laced simple finite-dimensional \(\mathfrak g\), does there exist a reduced decomposition, denoted by \(w_0\) in the supplied statement, for which
\[
S_{\mathrm{gm}}\cap D^q_{w_0}\ne\varnothing?
\]
\end{enumerate}
Here simply laced refers to a root system whose Dynkin diagram has only single edges, assuming the intended root-system setting. The record supplies neither the ground field and its characteristic nor the definitions of global monomial degrees, \(S_{\mathrm{gm}}\), and \(D^q_{w_0}\). It also does not explain \(q\), the ambient space of the intersection, or the choices on which these sets depend. The notation \(w_0\) does not distinguish a Weyl-group element from a chosen reduced expression, and the element being decomposed is not explicitly identified. These are unresolved qualifications: the displayed questions preserve the supplied targets but do not constitute a fully specified standalone conjecture without the missing source definitions.
\subsection{Short English statement}
The problem asks whether every simple finite-dimensional Lie algebra admits a global monomial degree and whether, in every simply laced case, such a degree can also belong to the set \(D^q_{w_0}\) for some reduced decomposition. The supplied record leaves the defining parameter spaces and ground-field assumptions unspecified.
\subsection{Sources}
\begin{itemize}
\item Supplied local UnsolvedMath registry record, permanent ID \(30003068\), exact source problem code \texttt{catalogue entry}, titled \emph{Existence of Global Monomial Orders for PBW Bases}. The record assigns the problem to 2016 and supplies no individual proposer.
\item \emph{Mini-Workshop: PBW Structures in Representation Theory} (2016), Oberwolfach Reports, pp.~615--651, DOI \texttt{10.4171/owr/2016/13}. Supplied source and literature link: \url{https://doi.org/10.4171/owr/2016/13}. The record describes the report as documenting workshop discussions and open problems.
\item Supplied literature reference labeled \emph{essential bases and toric degenerations}: \url{https://arxiv.org/abs/1510.02295}. No authors or more detailed bibliographic attribution are supplied.
\item Literature assessment dated exactly \texttt{2026-08-21}: the supplied curation labels the problem partially solved and states that nonemptiness is proved for several Lie types listed in the report, while the assertion for every simple Lie algebra and the stronger simply laced assertion remain unproved in general. This is supplied dated curation derived from the OWR material, not an independently verified literature assessment; the record requests independent verification of time-sensitive status and cited claims.
\end{itemize}
\end{document}
